\documentclass[10pt,a4paper]{amsart}
\usepackage[margin=19mm]{geometry}
\usepackage{amsmath,amssymb,mathtools}
\usepackage[scr=rsfs,cal=euler]{mathalfa}
\usepackage{xfrac}
\usepackage[unicode,hidelinks,bookmarks=false]{hyperref}
\newcommand{\ie}{i.e.\ }
\newcommand{\cf}{cf.\ }
\newcommand{\resp}{respectively\ }
\newcommand{\Subsection}{\subsection}
\providecommand{\nobreakdash}{\nobreak\hbox{-}}
\setlength{\emergencystretch}{2em}
\multlinegap0pt
\usepackage{bm}
\usepackage{bbm}
\usepackage{dsfont}
\usepackage{xspace}
\usepackage{cancel}
\usepackage{mathtools}
\usepackage{amsthm}
\makeatletter
\@ifundefined{thm}{\newtheorem{thm}{Theorem}}{}
\@ifundefined{lemma}{\newtheorem{lemma}[thm]{Lemma}}{}
\makeatother
\DeclareMathOperator{\Hom}{Hom}
\DeclareMathOperator{\id}{id}
\newcommand{\G}{\mathcal{G}}
\newcommand{\K}{\mathcal{K}}
\newcommand{\T}{\mathcal{T}}
\newcommand{\ccinfty}{C_c^\infty}
\title[Equivariant periodic cyclic homology for ample groupoids]{Equivariant periodic cyclic homology for ample groupoids}
\author{Francesco Pagliuca, Christian Voigt}
\date{Source: arXiv:2506.06503v2 (CC BY 4.0)}
\begin{document}
\maketitle

\noindent\emph{Source and licence.} Equivariant periodic cyclic homology for ample groupoids. Version arXiv:2506.06503v2 (CC BY 4.0), distributed under the Creative Commons Attribution 4.0 International licence, as recorded in the arXiv licence metadata for this exact version and shown on the abstract page https://arxiv.org/abs/2506.06503v2. Full rights notice: licenses/arxiv-2506.06503-CC-BY-4.0.txt.\par\smallskip

\noindent\textit{Editorial scope.} Counted unit: the Thm whose source label is \texttt{StabLemma} in the article (source0.tex lines 1577--1583, source label \texttt{StabLemma}), delivered together with the article's own complete original proof of it (source0.tex lines 1585--1653, one \texttt{proof} environment of 69 lines). No mathematics is written, repaired or re-derived: statement and proof are the article's own text line for line. Required context retained uncounted: thm: homotopy invariance (lines 1504--1507); twistedtrace (lines 1552--1558). Every definition, display and label of those lines is retained unchanged. Omitted from this package and not redistributed: 1--1503, 1508--1551, 1559--1576, 1584--1584, 1654--2132. Nothing inside a retained excerpt is omitted or altered: every retained body is delivered exactly as the article's lines are written, so no omission receipt of any kind accompanies this package. Citations: the retained excerpts carry no citation command at all, so no bibliography entry of the article is transcribed and no reference is redistributed. Reference adaptation: none was needed and none was made; every reference inside the delivered unit resolves inside it, so no unresolved reference or citation is produced. Printed slips kept literal: no printed word, symbol, hypothesis or number of the counted statement or of its proof was altered. Layout and class: the article's own class is \texttt{amsart} with packages including amsmath, graphicx, comment, stackrel, amssymb, MnSymbol, color, amscd, tikz, tikz-cd; the wrapper is the lane's pinned \texttt{amsart} editorial wrapper at 10pt on A4, which restates the theorem-like environments the retained text needs and adds the article's own macro definitions that the retained text needs (\texttt{\textbackslash G}, \texttt{\textbackslash Hom}, \texttt{\textbackslash K}, \texttt{\textbackslash T}, \texttt{\textbackslash ccinfty}, \texttt{\textbackslash id}), with extra wrapper packages (bm, bbm, dsfont, xspace, cancel, mathtools). The delivered numbering is the unit's own. Assets: no retained span contains a figure environment, a graphics-inclusion command, a PDF-page inclusion, a table or a TikZ picture; no image file is copied into this package, the package declares no asset dependency and no image file is redistributed with it. Licence: version arXiv:2506.06503v2 of this article is distributed under the Creative Commons Attribution 4.0 International licence, as recorded in the arXiv licence metadata for this exact version and shown on the abstract page https://arxiv.org/abs/2506.06503v2. A full rights notice is delivered at licenses/arxiv-2506.06503-CC-BY-4.0.txt. Characters: every retained excerpt is pure ASCII, so the delivered engine, XeLaTeX with its default Latin Modern text font, reports no missing character.
\begin{thm}[Homotopy invariance] \label{thm: homotopy invariance}
Let $ A $ and $ B $ be pro-$ \G $-algebras and let $ \Phi: A \rightarrow B[0,1] $ be a $ \G $-equivariant homotopy. Then the elements $ \left[\Phi_0\right] $ and $ \left[\Phi_1\right] $ in $ HP_0^\G(A,B) $ are equal. More generally, if $ A $ is a quasifree pro-$ \G $-algebra then the elements $ \left[\Phi_0\right] $ and $ \left[\Phi_1\right] $ 
in $ H_0(\Hom_{A(\G)}(X_\G(A), X_\G(B))) $ are equal.    
\end{thm}

\begin{lemma} \label{twistedtrace}
The twisted trace map introduced above satisfies
$$
ttr(\chi_U \otimes L_0 L_1) = ttr(\chi_U \otimes (\chi_{U^{-1}}\cdot L_1) L_0)
$$
for any compact open bisection $ U \subseteq \G $ and $ L_0, L_1 \in \K(E) $.  
\end{lemma}

\begin{thm} \label{StabLemma} 
Let $ A $ be a pro-$ \G $-algebra, and let $ E $ be a $ \G $-module equipped with a $ \G $-admissible $ \G $-equivariant bilinear pairing. Then the class 
$$
[\iota_A] \in H_0 \Hom_{A(\G)}(X_\G(\T A), X_\G(\T(A \otimes_{\ccinfty(\G^{(0)})} \K(E)))) 
$$
is invertible.
\end{thm}

\begin{proof}
We have to find an inverse for $ [\iota_A] $. First observe that the canonical $ \G $-equivariant linear map $ A \otimes_{\ccinfty(\G^{(0)})} \K(E) \rightarrow 
\T A \otimes_{\ccinfty(\G^{(0)})} \K(E) $ is a lonilcur and hence induces 
a $ \G $-equivariant homomorphism $ \lambda_A: \T(A \otimes_{\ccinfty(\G^{(0)})}
\K(E)) \rightarrow \T A \otimes_{\ccinfty(\G^{(0)})} \K(E) $.
Define $ tr_A: X_\G(\T A \otimes_{\ccinfty(\G^{(0)})} \K(E))
\rightarrow X_\G(\T A) $ by
\begin{equation*}
tr_A(f \otimes x \otimes L) = ttr(f \otimes L) \otimes x
\end{equation*}
and
\begin{align*}
tr_A(f \otimes (x_0 \otimes L_0) d(x_1 \otimes L_1)) &= 
ttr(f \otimes L_0 L_1) \otimes x_0 dx_1 \\
tr_A(f \otimes d(x_1 \otimes L_1)) &= 
ttr(f \otimes L_1) \otimes dx_1.
\end{align*}
Here, we use the twisted trace $ ttr $, see Lemma \ref{twistedtrace}. By construction, $ tr_A $ is a map of $ \G $-anti-Yetter-Drinfeld modules. 
We have
\begin{align*}
tr_A d_\G(f \otimes x \otimes L) &= tr_A(f \otimes d(x \otimes L)) \\
&= ttr(f \otimes L) \otimes dx \\
&= d_\G(ttr(f \otimes L) \otimes x) \\
&= d_\G tr_A(f \otimes x \otimes L), 
\end{align*}
and for a compact open bisection $ U \subseteq \G $ we calculate 
\begin{align*}
b_\G tr_A&(\chi_U \otimes (x_0 \otimes L_0) d(x_1\otimes L_1)) = b_\G(ttr(\chi_U \otimes L_0 L_1) \otimes x_0dx_1) \\
&= ttr(\chi_U \otimes L_0 L_1) \otimes (x_0x_1 - (\chi_{U^{-1}}\cdot x_1)x_0) \\
&= ttr(\chi_U \otimes L_0 L_1) \otimes x_0x_1 
- ttr(\chi_U \otimes (\chi_{U^{-1}}\cdot L_1) L_0) \otimes (\chi_{U^{-1}}\cdot x_1)x_0 \\ 
&= tr_A(\chi_U \otimes (x_0x_1\otimes L_0 L_1) - \chi_U \otimes (\chi_{U^{-1}} \cdot x_1)x_0 \otimes (\chi_{U^{-1}}\cdot L_1)L_0) \\
&= tr_A b_\G(\chi_U \otimes (x_0\otimes L_0) d(x_1\otimes L_1)), 
\end{align*}
using the twisted trace property from Lemma \ref{twistedtrace}. Similarly one checks
$$
b_\G tr_A(\chi_U \otimes d(x_1\otimes L_1)) = tr_A b_\G(\chi_U \otimes d(x_1\otimes L_1)). 
$$
It follows that $ tr_A $ is a chain map of paracomplexes. 

We define $ \tau_A = tr_A X_\G(\lambda_A) $ and claim that $ [\tau_A] $ is an inverse for $ [\iota_A] $. From the definitions one immediately computes 
$ [\iota_A] \cdot [\tau_A] = \id $. 
It thus remains to show that $ [\tau_A] \cdot
[\iota_A] = \id $. Consider the $ \G $-equivariant homomorphisms $ i_j: A \otimes_{\ccinfty(\G^{(0)})} \K(E) \rightarrow A \otimes_{\ccinfty(\G^{(0)})} \K(E) \otimes_{\ccinfty(\G^{(0)})} \K(E) $ for $ j = 0,1 $ given by 
\begin{align*}
i_0 = \id \otimes \iota, \quad 
i_1 = (\id \otimes \sigma) i_0. 
\end{align*}
Here we use the canonical identification 
$$ 
A \otimes_{\ccinfty(\G^{(0)})} \K(E) \cong A \otimes_{\ccinfty(\G^{(0)})} \K(E) \otimes_{\ccinfty(\G^{(0)})} \ccinfty(\G^{(0)}) 
$$ 
in the definition of $ i_0 $, and the tensor flip automorphism $ \sigma $ of $ \K(E) \otimes_{\ccinfty(\G^{(0)})} \K(E) $ given by $ \sigma(L_1 \otimes L_2) = L_2 \otimes L_1 $ in the definition of $ i_1 $. 

We calculate $ [i_0] \cdot [\tau_{A\otimes \K(E)}] = \id  $ 
and $ [i_1] \cdot [\tau_{A \otimes \K(E)}] =
[\tau_A] \cdot [\iota_A] $. Let us show that the maps $ i_0 $ and $ i_1 $ are $ \G $-equivariantly homotopic. To this end observe that $ \K(E) \otimes_{\ccinfty(\G^{(0)})} \K(E) \cong \K(E \otimes_{\ccinfty(\G^{(0)})} E)$ as $ \G $-algebras and denote by $ \Sigma $ the flip automorphism of $ E \otimes_{\ccinfty(\G^{(0)})} E $ given by $ \Sigma (e \otimes f) = f \otimes e $. For $ t \in [0,1] $ we then obtain a $ \G $-equivariant linear endomorphism $ \Sigma_t $ of $ E \otimes_{\ccinfty(\G^{(0)})} E $ given by
$$ 
\Sigma_t = \cos(\pi t/2) \id + i \sin(\pi t/2) \Sigma,
$$
and we note that $ \Sigma_t $ is invertible with inverse $ \Sigma_t^{-1} = \cos(\pi t/2) \id - i \sin(\pi t/2) \Sigma $. Conjugation with $ \Sigma_t $ defines $ \G $-equivariant algebra automorphism $ \sigma_t $ of 
$$ 
\K(E \otimes_{\ccinfty(\G^{(0)})} E) = (E \otimes_{\ccinfty(\G^{(0)})} E) \otimes_{\ccinfty(\G^{(0)})} (E \otimes_{\ccinfty(\G^{(0)})} E). 
$$ 
The family $ (\sigma_t)_{t \in [0,1]} $ depends smoothly on $ t $, and by construction we have 
$ \sigma_0 = \id $ and $ \sigma_1 = \sigma $. Now define $ h_t: A \otimes_{\ccinfty(\G^{(0)})} \K(E) \rightarrow 
A \otimes_{\ccinfty(\G^{(0)})} \K(E) \otimes_{\ccinfty(\G^{(0)})} \K(E) $ by $ h_t = (\id \otimes \sigma_t) i_0 $ for $ t \in [0,1] $. Then each $ h_t $ is a $ \G $-equivariant algebra homomorphism, and by construction
$ h_j = i_j $ for $ j = 0,1 $. Since the family $ (h_t)_{t \in [0,1]} $ depends again smoothly on $ t $ we have thus constructed a $ \G $-equivariant smooth homotopy between $ i_0 $ and $ i_1 $. According to Theorem \ref{thm: homotopy invariance} we obtain $ [i_0] = [i_1] $, and hence $ [\tau_A] \cdot [\iota_A] = \id $ as required. 
\end{proof}

\end{document}
